
\documentclass[11pt]{article}
\usepackage[margin=1in]{geometry}
\usepackage{amsmath,amssymb,booktabs,hyperref,graphicx,parskip,enumitem}
\usepackage{xcolor}
\hypersetup{colorlinks=true,linkcolor=blue!50!black,urlcolor=blue!50!black}
\setlist{nosep}

\title{Replication Report: (title not recorded)\\[2pt]\large OSTI 2396968 \textbar{} Coverage 9/10, Agreement 6/10}
\author{Rick Stevens \and Ollie (AI Assistant)\\\small Argonne National Laboratory}
\date{2026-04-27}

\begin{document}\maketitle

\section{Source paper}
\begin{itemize}
\item \textbf{Title:} (title not recorded)
\item \textbf{Authors:} Not recorded
\item \textbf{Year:} n/a
\item \textbf{Domain:} Not recorded
\item \textbf{Identifier:} OSTI 2396968
\end{itemize}


\section{Replication objective}
The central claim of the paper concerns not recorded. Our objective was to reproduce the headline numerical/qualitative results within the constraints of the AI-driven Replication Project (24--72\,h, commodity GPU/CPU compute, no author contact). A replication plan is archived at \texttt{2396968-Latent-Stochastic-Differential-Equations-for-Modeling/replication\_plan.pdf}.

\section{Methodology}
We followed the standard Replication Project workflow: ingest the source PDF, extract method and data pointers, draft a replication plan, attempt the smallest end-to-end pipeline that produces a comparable number, then scale up where compute and data permit. Tools used in this replication are catalogued in the meta-paper's computational tools inventory (see \texttt{COMPUTATIONAL\_TOOLS\_INVENTORY.md}).

This paper went through the Tier-Lift pipeline; supplementary notes at \texttt{.openclaw/workspace/24h-progress/tier-lift-v2/2396968.md}.

\subsection*{Paper-directory README excerpt}
\small
\par\medskip\textbf{Latent Stochastic Differential Equations for Modeling Quasar Variability and Inferring Black Hole Properties}\par
\par$\bullet$\ **OSTI ID:** 2396968
\par$\bullet$\ **Rank:** \#27
\par$\bullet$\ **Replication Score:** 9/10 (paper); **v1 (1-band, deprecated) 4-6/10**; **v2 (paper-faithful, 6-band, A100) 7/10**
\par$\bullet$\ **v2 report:** [`replication/v2\_faithful/report/report.pdf`](replication/v2\_faithful/report/report.pdf)
\par$\bullet$\ **v2 training:** 18.08 h on uicgpu A100, 8 epochs, 20k curves, 917,594 params, full 6-band + Sim5 GR + 9-param Cholesky head
\par$\bullet$\ **v2 headline:** LC RMSE 0.1091 mag (paper 0.0959, 1.14×); 1σ/2σ/3σ coverage 73.3/94.5/99.0\% (paper \textasciitilde{}70/95/99.7\%); per-parameter ranking matches paper exactly
\par$\bullet$\ **Open-Source Tools:** Yes (PyTorch + torchsde)
\par$\bullet$\ **Code Repository:** No (paper); this directory contains a from-scratch simplified replication

\par\medskip\textbf{\# What was replicated}\par

A single-band simplification of Fagin et al. (2024): encoder → latent Itô SDE (learned prior drift,
diagonal diffusion) → decoder, trained with reconstruction NLL plus Girsanov KL via
`torchsde.sdeint(..., logqp=True)`. We skip the 6-band LSST pipeline, the physical accretion-disk
transfer functions, and the 9-parameter black-hole inference head.

\par\medskip\textbf{\# Quantitative summary}\par

| Model | RMSE (mag) | MAE (mag) | NLL |
|---|---|---|---|
| Paper latent-SDE (6-band) | 0.0959 | 0.0695 | −1.14 |
| Paper GPR (Matern-½, 6-band) | 0.0978 | 0.0711 | −1.01 |
| **Ours latent-SDE (1-band, 512 train)** | 0.198 | 0.149 | +0.49 |
| **Ours GPR (Matern-½, LOO, 1-band)** | 0.048 | 0.034 | −1.85 |

The gap to the paper is driven by (i) single-band setup — the paper's edge over GPR largely comes
from sharing information across six bands, (ii) a \textasciitilde{}2000× smaller training set, (iii) an \textasciitilde{}17×
smaller model, and (iv) \textasciitilde{}3000× less training compute (160 s CPU vs \textasciitilde{}6 weeks V100).

\par\medskip\textbf{\# Layout}\par

[code block omitted]

\par\medskip\textbf{\# How to reproduce (CPU, \textasciitilde{}3 min total)}\par

[code block omitted]

\par\medskip\textbf{\# Status}\par
\par$\bullet$\ [x] Paper reviewed
\par$\bullet$\ [x] Code/tools identified (PyTorch + torchsde)
\par$\bullet$\ [x] Code implemented (from scratch — no public repo)
\par$\bullet$\ [x] Simplified results reproduced (single-band DRW)
\par$\bullet$\ [ ] Full-scale results validated against paper (out of scope: 6 weeks V100)

\par\medskip\textbf{\# Notes / Caveats}\par
\par$\bullet$\ `torchsde.sdeint(..., logqp=True)` is exactly the Li-et-al. 2020 / Fagin-et-al. formulation for
  the Girsanov KL between posterior and prior SDE paths.

[\textbackslash{}textit\{README truncated\}]
\normalsize

The replication was driven by an OpenClaw agent loop (Claude Opus / GPT-5 class models) issuing shell, Python, and (where applicable) GPU jobs. Compute usage and wall-time for individual papers are summarised in the meta-paper's resource table; not separately recorded for this paper unless noted in the Tier-Lift artefact. Sub-agents were spawned for replication-plan generation and (for papers in batches 1--4) for tier-lift extension runs.

\section{What was reproduced}
\begin{itemize}
\item Not recorded.
\end{itemize}

\section{What was missing or incomplete}
\begin{itemize}
\item Not recorded.
\end{itemize}
The dominant blocker categories for this paper, mapped onto the meta-paper's 9-category friction taxonomy (\S4), are listed in \S\ref{sec:friction} below.

\section{Quantitative agreement}
Per-quantity numerical comparisons are not separately tabulated in the structured evaluation record for this paper beyond the agreement rationale below; where the rationale references specific numbers, they are reproduced verbatim. A full numerical comparison table is recorded in the per-replication artefacts (see \S\ref{sec:repro}). For papers below 8/8, the agreement rationale identifies the specific quantities that diverge.

\paragraph{Agreement rationale.} Not recorded.

\section{Score and friction-point tagging}\label{sec:friction}
\begin{itemize}
\item \textbf{Coverage:} 9/10 --- Not recorded.
\item \textbf{Agreement:} 6/10 --- Not recorded.
\end{itemize}

\noindent\textbf{Friction points encountered (taxonomy tags):}
\begin{itemize}
\item None recorded.
\end{itemize}

\section{Follow-on questions}
\begin{itemize}
\item Not recorded.
\end{itemize}

\section{Reproducibility pointers}\label{sec:repro}
\begin{itemize}
\item \textbf{Paper directory:} \texttt{2396968-Latent-Stochastic-Differential-Equations-for-Modeling}
\item \textbf{Replication artefacts:}
\begin{itemize}
\item \texttt{/Users/stevens/Dropbox/REPLICATE-PROJECT/2396968-Latent-Stochastic-Differential-Equations-for-Modeling/replication}
\item \texttt{/Users/stevens/Dropbox/REPLICATE-PROJECT/2396968-Latent-Stochastic-Differential-Equations-for-Modeling/replication/v2\_faithful}
\item \texttt{/Users/stevens/Dropbox/REPLICATE-PROJECT/2396968-Latent-Stochastic-Differential-Equations-for-Modeling/replication/v1\_simplified}
\end{itemize}
\item \textbf{Replication-dir hint from JSONL:} \texttt{/Users/stevens/Dropbox/REPLICATE-PROJECT/2396968-Latent-Stochastic-Differential-Equations-for-Modeling}
\item \textbf{Status:} tier\_lift\_v2\_2026\_04\_27
\end{itemize}

To re-run, change directory into the paper directory above and consult its \texttt{README.md} for the entry-point script. Most replications follow the convention \texttt{replication/run\_*.sh} or \texttt{replication/<subdir>/run.py}.

\end{document}
